% =============================================================================
% 毕业论文LaTeX模板
% FormalMath项目学术写作支持系统
% 适用：本科/硕士毕业论文（50+页）
% =============================================================================

\documentclass[12pt,a4paper]{report}

% ==================== 宏包引入 ====================
% 中文支持
\usepackage{ctex}

% 页面设置
\usepackage{geometry}
\geometry{left=3cm,right=3cm,top=2.5cm,bottom=2.5cm}

% 数学宏包
\usepackage{amsmath,amssymb,amsthm,mathtools}

% 定理环境
\theoremstyle{definition}
\newtheorem{theorem}{定理}[chapter]
\newtheorem{lemma}[theorem]{引理}
\newtheorem{proposition}[theorem]{命题}
\newtheorem{corollary}[theorem]{推论}
\newtheorem{definition}[theorem]{定义}
\newtheorem{example}[theorem]{例子}
\newtheorem{remark}[theorem]{注记}
\newtheorem{conjecture}[theorem]{猜想}

% 图表
\usepackage{graphicx}
\usepackage{tikz}
\usepackage{tikz-cd}
\usepackage{pgfplots}
\pgfplotsset{compat=1.18}
\usetikzlibrary{arrows,shapes,positioning,calc}

% 参考文献
\usepackage[numbers,sort&compress]{natbib}
\bibliographystyle{amsplain}

% 超链接
\usepackage{hyperref}
\hypersetup{
    colorlinks=true,
    linkcolor=blue,
    citecolor=blue,
    urlcolor=blue,
    bookmarksnumbered=true,
    pdftitle={毕业论文标题},
    pdfauthor={作者姓名}
}

% 页眉页脚
\usepackage{fancyhdr}
\pagestyle{fancy}
\fancyhf{}
\fancyhead[LE,RO]{\thepage}
\fancyhead[LO]{\small\nouppercase{\leftmark}}
\fancyhead[RE]{\small\nouppercase{\rightmark}}
\renewcommand{\headrulewidth}{0.4pt}

% 行间距
\usepackage{setspace}
\onehalfspacing

% 列表
\usepackage{enumitem}

% 算法
\usepackage{algorithm}
\usepackage{algorithmic}

% 代码（如需要）
\usepackage{listings}

% 符号表
\usepackage{nomencl}
\makenomenclature

% 索引（可选）
\usepackage{makeidx}
\makeindex

% ==================== 文档信息 ====================
\title{\Huge\textbf{毕业论文标题}\\[1em]
\Large 副标题（如有）}
\author{作者姓名}
\date{\today}

% ==================== 正文开始 ====================
\begin{document}

% ==================== 封面 ====================
\begin{titlepage}
    \centering
    \vspace*{2cm}

    {\Huge\bfseries 学校名称}

    \vspace{1cm}

    {\LARGE 院系名称}

    \vspace{2cm}

    \rule{\textwidth}{1pt}

    \vspace{0.5cm}

    {\Huge\bfseries 毕业论文}

    \vspace{0.5cm}

    \rule{\textwidth}{1pt}

    \vspace{2cm}

    {\LARGE\bfseries 论文标题}

    \vspace{0.5cm}

    {\Large 副标题（如有）}

    \vspace{3cm}

    \begin{tabular}{rl}
        \textbf{作者：} & 姓名 \\[0.5em]
        \textbf{学号：} & 12345678 \\[0.5em]
        \textbf{专业：} & 数学与应用数学 \\[0.5em]
        \textbf{导师：} & 导师姓名 教授 \\[0.5em]
        \textbf{日期：} & \today
    \end{tabular}

    \vfill
\end{titlepage}

% ==================== 摘要 ====================
\chapter*{摘要}
\addcontentsline{toc}{chapter}{摘要}

这是中文摘要。应该包含：
\begin{itemize}
    \item 研究背景和动机
    \item 研究问题
    \item 主要方法
    \item 主要结果
    \item 研究意义
\end{itemize}

\vspace{1em}
\noindent\textbf{关键词：} 关键词1；关键词2；关键词3

\newpage

% ==================== Abstract ====================
\chapter*{Abstract}
\addcontentsline{toc}{chapter}{Abstract}

This is the English abstract. It should mirror the content of the Chinese abstract.

\vspace{1em}
\noindent\textbf{Keywords:} keyword1; keyword2; keyword3

\newpage

% ==================== 致谢 ====================
\chapter*{致谢}
\addcontentsline{toc}{chapter}{致谢}

感谢导师的悉心指导...感谢家人的支持...感谢同学的帮助...

\newpage

% ==================== 目录 ====================
\tableofcontents
\newpage

% ==================== 符号表 ====================
\chapter*{符号表}
\addcontentsline{toc}{chapter}{符号表}

\begin{tabular}{ll}
    $\mathbb{N}$ & 自然数集 \\
    $\mathbb{Z}$ & 整数集 \\
    $\mathbb{Q}$ & 有理数集 \\
    $\mathbb{R}$ & 实数集 \\
    $\mathbb{C}$ & 复数集 \\
    $\mathbb{F}_q$ & 有$q$个元素的有限域 \\
    $GL_n(F)$ & 域$F$上的一般线性群 \\
    $SL_n(F)$ & 域$F$上的特殊线性群
\end{tabular}

\newpage

% ==================== 第1章 引言 ====================
\chapter{引言}\label{ch:intro}

\section{研究背景}

介绍研究主题的背景，从一般到具体。

\section{问题陈述}

明确陈述本文要研究的具体问题。

\section{文献综述}

\subsection{早期工作}

回顾相关领域的历史发展。

\subsection{近期进展}

介绍最新的研究成果。

\section{本文贡献}

明确列出本文的主要贡献：
\begin{enumerate}
    \item 贡献1
    \item 贡献2
    \item 贡献3
\end{enumerate}

\section{论文结构}

本文结构如下：
\begin{itemize}
    \item 第\ref{ch:prelim}章介绍必要的预备知识
    \item 第\ref{ch:main}章陈述并证明主要结果
    \item 第\ref{ch:applications}章讨论应用和例子
    \item 第\ref{ch:conclusion}章总结全文并展望未来
\end{itemize}

% ==================== 第2章 预备知识 ====================
\chapter{预备知识}\label{ch:prelim}

本章回顾本文所需的定义和已知结果。

\section{基本概念}

\begin{definition}\label{def:fundamental}
基本概念的定义。
\end{definition}

\begin{example}
说明基本概念的简单例子。
\end{example}

\section{主要工具}

\begin{theorem}[经典结果]\label{thm:classical}
重要的经典定理。
\end{theorem}

\begin{proof}
证明或参考文献。
\end{proof}

% ==================== 第3章 主要结果 ====================
\chapter{主要结果}\label{ch:main}

本章陈述并证明本文的主要结果。

\section{主要定理}

\begin{theorem}[主要结果]\label{thm:main}
本文的主要定理。
\end{theorem}

\begin{proof}
详细的证明过程...
\end{proof}

\section{推论与应用}

\begin{corollary}\label{cor:first}
第一个重要推论。
\end{corollary}

% ==================== 第4章 应用 ====================
\chapter{应用与例子}\label{ch:applications}

\section{具体例子}

通过具体例子说明理论的应用。

\section{数值结果（如适用）}

如果涉及计算，展示数值结果。

% ==================== 第5章 结论 ====================
\chapter{结论与展望}\label{ch:conclusion}

\section{工作总结}

总结本文的主要成果。

\section{局限性与未来工作}

讨论本文的局限性，并提出未来研究的方向。

% ==================== 参考文献 ====================
\newpage
\bibliography{references}

% ==================== 附录 ====================
\appendix
\chapter{补充证明}

这里可以放置详细的技术性证明。

\chapter{计算细节}

如果需要，放置详细的计算过程。

\end{document}
% =============================================================================
% 使用说明：
% 1. 使用XeLaTeX编译
% 2. 需要创建references.bib文件
% 3. 根据学校要求调整封面格式
% 4. 编译顺序：XeLaTeX → BibTeX → XeLaTeX × 2
% =============================================================================
